\chapter{Componenta Hardware}
\label{ch:hardware}

\section{LattePanda Mu}

LattePanda Mu reprezintă un modul de calcul de înaltă performanță, cu dimensiuni comparabile cu cele ale unui card bancar, proiectat de compania DFRobot. Conceptul principal al acestui modul este de a facilita integrarea în produse hardware customizate, fiind pretabil atât pentru proiecte de tip hobby, cât și pentru aplicații industriale.

Seria actuală cuprinde modele cu diferite configurații de procesoare și memorie, acestea menținând intercompatibilitatea la nivel de pinout. Modulul este proiectat pentru a fi inserat într-o placă de bază (carrier board) folosind un slot standard SO-DIMM \ac{DDR}4. Această alegere constructivă este justificată de densitatea mare de pini, necesară pentru a expune numeroasele interfețe de comunicare către placa purtătoare.

Flexibilitatea platformei este sporită de posibilitatea modificării \ac{BIOS}-ului, permițând utilizatorului să configureze alocarea liniilor de date pentru a prioritiza, după caz, numărul de sloturi \ac{PCIe} sau porturile \ac{USB} disponibile. Specificațiile tehnice ale unității utilizate în prezenta lucrare sunt următoarele:




În ceea ce privește conectivitatea, configurația aleasă pentru acest sistem include:
\begin{itemize}
    \item 2 porturi \ac{USB} 3.2 Gen 2 (10 Gbps);
    \item 2 ieșiri video \ac{HDMI} 2.0;
    \item DisplayPort 1.4 prin interfața \ac{USB}-C;
    \item 7 linii \ac{PCIe} Gen 3.
\end{itemize}

\section{Conexiuni - Diagramă Bloc}

În această secțiune sunt descrise conexiunile și este explicată diagrama bloc din Figura ~\ref{fig:diagrama-bloc}.
Sistemul utilizează standardul \textit{\ac{USB} \ac{PD}}  pentru alimentare.
Această soluție a fost aleasă datorită costului redus și a disponibilității ridicate pe piață a surselor compatibile.
Fiind un standard universal și nu unul proprietar, acesta oferă flexibilitate maximă în alegerea alimentatorului și facilitează înlocuirea rapidă a acestuia în caz de necesitate.
Toate conexiunile sunt realizate intern, la nivelul plăcii de bază, pentru a reduce latențele și a crește fiabilitatea.


\begin{figure}[H]
  \centering
  \begin{subfigure}[b]{1\linewidth}
    \includegraphics[width=\linewidth]{ports-front.jpg}
     \caption{Port-urile din partea din față a PCB-ului}
  \end{subfigure}
  \begin{subfigure}[b]{1\linewidth}
    \includegraphics[width=\linewidth]{ports-back.jpg}
    \caption{Port-urile din partea din spate a PCB-ului}
  \end{subfigure}
  \begin{subfigure}[b]{0.7\linewidth}
    \includegraphics[width=\linewidth]{pcb-bottom.jpg}
    \caption{Port-urile de pe spatele PCB-ului}
  \end{subfigure}
  \caption{Port-urile PCB-ului}
  \label{fig:ports}
\end{figure}



\begin{figure}[h!]
    \makebox[\textwidth][c]{\includegraphics[width=1.15\textwidth]{MiniLatte_block.png}}
    \caption{Diagramă bloc a sistemului MiniLatte. Prezintă interconexiunile între modulul LattePanda Mu și componentele periferice.}
    \label{fig:diagrama-bloc}
\end{figure}

Pentru extinderea capacităților de conectare a perifericelor, au fost integrate două hub-uri \ac{USB} care pun la dispoziție 4 porturi \ac{USB} 3.2 Gen 2. În ceea ce privește conectivitatea la rețea, sistemul oferă două posibilități distincte, ambele bazate pe magistrala \ac{PCIe} Gen 3:
\begin{enumerate}
    \item \textbf{Wireless:} prin intermediul unui modul \ac{Wi-Fi} instalat în slotul dedicat de pe placa de bază;
    \item \textbf{Ethernet:} utilizând controlerul \textit{Realtek RTL8125}, care suportă viteze de transfer de până la 2.5 Gbps via \ac{PCIe} Gen 3.
\end{enumerate}

Capacitatea de stocare poate fi extinsă considerabil prin cele trei sloturi destinate unităților \ac{SSD} M.2 \ac{NVMe}. Suplimentar, la unul dintre hub-urile \ac{USB} este conectat un cititor de carduri dual \ac{SD}, utilizând lățimea de bandă oferită de interfața \ac{USB} 3.2 Gen 2.

Pe partea de procesare audio, sistemul utilizează un cip \textit{ALC269Q} ce comunică prin protocolul \textit{\ac{HDAudio}}. Semnalul video este distribuit prin trei porturi digitale: două \ac{HDMI} 2.0 și un DisplayPort 1.4 (via \ac{USB}-C).



\section{Design}

Placa de bază pentru acest modul a fost proiectată folosind Kicad 9, folosind ca bază documentația dată de creatorii modulului, dar și țindând cont de regulile de proiectare atât pe partea de diagrame electrice dar și pe partea de design de circuit integrat imprimat (\ac{PCB}).

\subsection{Consum de curent}
Din cauza conectivității bogate, un lucru important in acest design a fost proiectarea in jurul consumuli de curent atât a modulului LattePanda Mu dar și componentelor de pe placa și a diverserlor dispozitive conectate.
Cum am menționat anterior, modulul este alimentat printr-un port \ac{USB} Type C prin protocolul \ac{USB} \ac{PD}, am folosit un chip HUSB238 pentru a comunica prin acest procol cu incărcătorul conectat și pentru a cere 20v 3,25A (65W) care este o valoare standard pentru acest protocol.
Aceasta tensiune este coborâtă la 17,5V pentru modulul LattePanda Mu folosind regulatorul TPS54561, consumul modulului fiind de 10-15W.
De asemenea tensiunea mai este coborâtă la 3,3V respectiv 5V folosind chip-ul TPS51275.
Planele de alimentare sunt foarte late, fiind nevoie de zone foarte late pentru a conduce curentul necesar. Acestea sunt 7-7,5mm dupa iesirea din regulatorul de tensiunea si dupa se ingusteaza din ce in ce mai tare, dupa ce trec de componentele cu consum mai mare. Pentru tensiunea de 3,3v acestea sunt cele 3 \ac{SD}-uri iar pentru cea de 5v sunt porturile \ac{USB} care pot consuma până la 0,5A pentru porturile \ac{USB} 2.0 și 0,9A pentru Porturile \ac{USB} 3.0 și Type C, acesea fiind monitorizate constant de mai multe chip-uri TPS2561  care pot sa întrerupă curent către portul \ac{USB} dacă detecteaza un consum prea mare de curent.  

\begin{figure}[H]
    \centering
    \includegraphics[width=\linewidth]{Latime plan curent.png}
    \caption{Plan de curent cu o lățime de 7,3mm.}
    \label{fig:power-plane}
\end{figure}


\subsection{Magistrale Diferențiale}
Magistralele diferețiale de comunicare sunt principala metoda in care sunt transmise date la viteze mari, marea majoritate a protocoalelor de comunicare moderne cum ar fi \ac{HDMI}, \ac{USB}, \ac{PCIe} folosesc acest mod de transmitere de date.
In proiectare de \ac{PCB}-uri, magistralele diferețiale reprezintă două linii paralele cu o lungime egală și cu o distanță foarte mică între ele, exista o toleranță între diferența de lungime a acestor linii (Intra-pair skew)  , acesta este de exemplu 0.125mm pentru USB 3.0 \cite{Layout_Guidelines}. Unele protocoale precum HDMI si DisplayPort care folosesc mai multe linii pentru output video trimit date sincron pe toate liniile, iar toate perechiile de linii trebuie sa fie egale ca lungime (Inter-Pair Skew)

\begin{figure}[H]
  \centering
  \begin{subfigure}[b]{0.3\linewidth}
    \includegraphics[width=\linewidth]{Intra-pair-skew.png}
     \caption{Corecție a Intra Pair Skew pe o Magistrală diferețială PCIe}
  \end{subfigure}
  \hspace{1cm}
  \begin{subfigure}[b]{0.6\linewidth}
    \includegraphics[width=\linewidth]{Inter-pair-skew.png}
    \caption{Corecție a Inter Pair Skew pe cele 4 magistrale diferețiale ale unui port HDMI}
  \end{subfigure}
  
  \caption{Corecții ale Intra si Inter Pair Skew}
  \label{fig:Skew}

\end{figure}

\section{Matrice LED}

Mini PC-ul vine echipat cu o matrice LED compusa din 220 LED-uri adresabile RGB cu driver WS2812C, acestea au nevoie doar de un pin de input și sunt legate în serie, acestea formând o matrice 44x5.
Acestea sunt controlate de un Microcontroller ESP32-C3 care comunica cu modulul LattePanda prin usb, iar microcontrollerul fiind plasat direct pe Placa de bază.

\begin{figure}[H]
    \centering
    \includegraphics[width=\linewidth]{Latime plan curent.png}
    \caption{Matrice Led. TO ADD PHOTO}
    \label{fig:LED-Matrix}
\end{figure}


\section{Carcasa}
Pentru a monta cele două \ac{PCB}-uri am proiectat o carcasă folosind AutoCAD, aceasta este formată din 2 piese principale printate 3d folosinf filament \ac{PLA}, un schelet ce include partea din față, lateralele și partea superioară, apoi un capac pentru partea inferioară și spate. Cele 2 \ac{PCB}-uri sunt montate pe schelet, acesta incluzand și un panou de difuzie printat cu filament translucid \ac{PETG}.




%$$